\documentclass[11pt,letterpaper]{article}

\usepackage[T1]{fontenc}
\usepackage[utf8]{inputenc}
\usepackage{mathpazo}
\usepackage[scaled=0.92]{helvet}
\usepackage{microtype}
\usepackage{geometry}
\usepackage{setspace}
\usepackage{amsmath,amssymb,amsthm,mathtools}
\usepackage{booktabs}
\usepackage{array}
\usepackage{enumitem}
\usepackage[most]{tcolorbox}
\usepackage{xcolor}
\usepackage{fancyhdr}
\usepackage{hyperref}
\usepackage{bookmark}

\geometry{
  top=0.88in,
  bottom=0.88in,
  inner=0.92in,
  outer=0.92in,
  headheight=15pt,
  headsep=22pt,
  footskip=30pt
}

\definecolor{JGPBlue}{HTML}{2E74B5}
\definecolor{JGPDarkBlue}{HTML}{1F4D78}
\definecolor{JGPInk}{HTML}{16283A}
\definecolor{JGPMuted}{HTML}{5D6670}
\definecolor{JGPLight}{HTML}{F4F6F9}
\definecolor{JGPBorder}{HTML}{CAD6E2}

\hypersetup{
  colorlinks=true,
  linkcolor=JGPDarkBlue,
  citecolor=JGPDarkBlue,
  urlcolor=JGPBlue,
  pdftitle={Local Fidelity, Global Alias},
  pdfsubject={A Pipeline Proof of the Three-Dimensional Jacobian Counterexample},
  pdfauthor={Jon Poplett},
  pdfkeywords={Jacobian conjecture, polynomial map, counterexample, local fidelity, global alias, proof pipeline}
}
\urlstyle{same}

\setstretch{1.075}
\setlength{\parindent}{0pt}
\setlength{\parskip}{0.54em}
\setlist[itemize]{leftmargin=1.65em,itemsep=0.23em,topsep=0.35em}
\setlist[enumerate]{leftmargin=1.85em,itemsep=0.23em,topsep=0.35em}
\setcounter{secnumdepth}{2}

\newcommand{\C}{\mathbb{C}}
\newcommand{\id}{\operatorname{id}}
\newcommand{\UPSN}{\textsc{Universal Proof Selection Network}}

\newtheoremstyle{jgpplain}
  {0.85em}
  {0.65em}
  {\itshape}
  {}
  {\bfseries\color{JGPDarkBlue}}
  {.}
  {0.5em}
  {}
\theoremstyle{jgpplain}
\newtheorem{theorem}{Theorem}
\newtheorem{lemma}[theorem]{Lemma}

\newtcolorbox{mechanismbox}{
  enhanced,
  breakable,
  colback=JGPLight,
  colframe=JGPBlue,
  boxrule=0pt,
  borderline west={2.2pt}{0pt}{JGPBlue},
  arc=0pt,
  left=10pt,
  right=10pt,
  top=7pt,
  bottom=7pt
}

\newtcolorbox{routebox}{
  enhanced,
  breakable,
  colback=white,
  colframe=JGPBorder,
  boxrule=0.55pt,
  arc=1pt,
  left=10pt,
  right=10pt,
  top=7pt,
  bottom=7pt
}

\newtcolorbox{auditbox}{
  enhanced,
  breakable,
  colback=JGPLight,
  colframe=JGPBorder,
  boxrule=0.55pt,
  arc=1pt,
  left=10pt,
  right=10pt,
  top=7pt,
  bottom=7pt
}

\pagestyle{fancy}
\fancyhf{}
\fancyhead[L]{\sffamily\bfseries\footnotesize\color{JGPMuted}LOCAL FIDELITY, GLOBAL ALIAS}
\fancyfoot[R]{\sffamily\footnotesize\color{JGPMuted}JGPTECH RESEARCH NOTE\enspace$\cdot$\enspace Page \thepage}
\renewcommand{\headrulewidth}{0pt}
\renewcommand{\footrulewidth}{0pt}

\makeatletter
\renewcommand{\section}{\@startsection{section}{1}{0pt}%
  {-2.7ex plus -0.8ex minus -0.2ex}%
  {1.15ex plus 0.2ex}%
  {\normalfont\sffamily\Large\bfseries\color{JGPBlue}}}
\renewcommand{\subsection}{\@startsection{subsection}{2}{0pt}%
  {-2.1ex plus -0.6ex minus -0.2ex}%
  {0.85ex plus 0.2ex}%
  {\normalfont\sffamily\large\bfseries\color{JGPDarkBlue}}}
\makeatother

\begin{document}

\hypersetup{pageanchor=false}
\begin{titlepage}
\thispagestyle{empty}
\vspace*{0.38in}

{\sffamily\bfseries\fontsize{29}{33}\selectfont\color{JGPInk}
Local Fidelity,\\[2pt]Global Alias\par}

\vspace{0.22in}
{\color{JGPBlue}\rule{\textwidth}{2.2pt}}
\vspace{0.22in}

{\sffamily\fontsize{16}{21}\selectfont\color{JGPDarkBlue}
A Pipeline Proof of the Three-Dimensional\\
Jacobian Counterexample\par}

\vspace{0.42in}

{\sffamily\bfseries\large Jon Poplett\par}
\vspace{0.08in}
{\sffamily\color{JGPMuted}
JGPTech Research Note \(\cdot\) 24 July 2026\\
AI-assisted proof construction and exact symbolic audit\par}

\vspace{0.52in}

\begin{mechanismbox}
\textbf{\color{JGPDarkBlue}Central mechanism.}
A polynomial representation may preserve every infinitesimal degree of freedom
while erasing a discrete global state label. The Jacobian detects the first
fact. A collision detects the second.
\end{mechanismbox}

\vfill

{\sffamily\small\color{JGPMuted}
\textbf{Primary author and investigator:} Jon Poplett\\
\textbf{Verification:} Exact polynomial arithmetic over the rational numbers
\par}
\end{titlepage}
\hypersetup{pageanchor=true}
\setcounter{page}{1}

\section*{Abstract}
\addcontentsline{toc}{section}{Abstract}

We give a self-contained proof that the Jacobian conjecture is false in
dimension three. The explicit degree-seven map is already publicly known; the
contribution of this note is a proof architecture that derives the certificate
from a general counterexample-construction pipeline.

The full state is a \emph{marked factorization} \((L,Q)\) of a binary cubic.
The visible representation is the product \(LQ\), which forgets which of the
cubic's three roots was assigned to the linear factor. A resultant
normalization removes the continuous scaling gauge and makes multiplication
locally faithful, while the discrete three-way ambiguity survives globally. A
specially aligned affine slice then closes under a two-level residual
hierarchy. Promoting those residuals to coordinates produces an affine
three-space and, after linear coordinate changes, the compact counterexample.

Thus the proof is routed as follows.

\begin{routebox}
\centering\sffamily\bfseries\color{JGPDarkBlue}
marked state \(\longrightarrow\) visible product
\(\longrightarrow\) gauge normalization
\(\longrightarrow\) local fidelity\\[2pt]
\(\longrightarrow\) global alias
\(\longrightarrow\) boundary closure
\(\longrightarrow\) finite certificate
\end{routebox}

No numerical approximation is used.

\section{Theorem intake and proof route}

The Jacobian conjecture over \(\C\) asserts that every polynomial map
\[
F:\C^n\longrightarrow\C^n
\]
with nonzero constant Jacobian determinant has a polynomial inverse.

The claim to be proved is existential.

\begin{theorem}\label{thm:counterexample}
There exists a polynomial map \(F:\C^3\to\C^3\) whose Jacobian determinant is
the nonzero constant \(-2\), but which is not injective. Consequently, the
Jacobian conjecture is false in dimension three and in every dimension
\(n\geq 3\).
\end{theorem}

The proof obligations are exact:
\begin{enumerate}
  \item construct a polynomial map \(F:\C^3\to\C^3\);
  \item prove \(\det DF\equiv -2\);
  \item exhibit distinct points with the same image.
\end{enumerate}

The pipeline adds four construction obligations before the final certificate.

\begin{center}
\renewcommand{\arraystretch}{1.22}
\begin{tabular}{>{\raggedright\arraybackslash}p{0.34\textwidth}
                >{\raggedright\arraybackslash}p{0.57\textwidth}}
\toprule
\sffamily\bfseries Pipeline stage &
\sffamily\bfseries Mathematical obligation\\
\midrule
Full state and visible representation &
Identify what information the map retains and what it erases\\
Invariant-preserving constraint &
Remove trivial continuous nonuniqueness without destroying the target collision\\
Local control &
Prove the constrained representation is infinitesimally faithful\\
Boundary closure &
Find polynomial coordinates covering the degenerate boundary and the generic
region simultaneously\\
\bottomrule
\end{tabular}
\end{center}

\section{The marked-factor representation}

Let \(U,V\) be homogeneous variables. Write
\[
L(U,V)=aU+bV,
\qquad
Q(U,V)=cU^2+dUV+eV^2.
\]
Their product is the binary cubic
\[
LQ
=ac\,U^3+(ad+bc)U^2V+(ae+bd)UV^2+be\,V^3.
\]
Therefore multiplication defines the polynomial representation
\[
M:\C^5\longrightarrow\C^4,
\qquad
M(a,b,c,d,e)=(ac,\ ad+bc,\ ae+bd,\ be).
\]

The input remembers a marked decomposition: one root belongs to \(L\), while
the other two belong to \(Q\). The output remembers only the product.

If a cubic has three distinct linear factors,
\[
C=L_1L_2L_3,
\]
then the three marked states
\[
(L_1,L_2L_3),\qquad
(L_2,L_1L_3),\qquad
(L_3,L_1L_2)
\]
all have visible representation \(C\). This is the global alias we want, but
multiplication also has a trivial continuous ambiguity:
\[
(L,Q)\longmapsto(\lambda L,\lambda^{-1}Q),
\qquad \lambda\in\C^\times.
\]
We remove that gauge before testing whether any genuine collision remains.

\section{Gauge normalization and local fidelity}

Define the resultant
\[
R(L,Q)
=R(a,b,c,d,e)
=a^2e-abd+b^2c.
\]
It vanishes exactly when \(L\) and \(Q\) share a projective root. Under the
gauge transformation above,
\[
R(\lambda L,\lambda^{-1}Q)=\lambda R(L,Q).
\]
Hence every state with \(R\neq0\) has a unique gauge representative satisfying
\[
R=1.
\]
Let
\[
Y=\{(a,b,c,d,e)\in\C^5:R(a,b,c,d,e)=1\}.
\]

\begin{lemma}[Multiplication is locally faithful on the normalized state space]
\label{lem:local-fidelity}
For every \(p\in Y\), the differential
\[
d(M|_Y)_p:T_pY\longrightarrow\C^4
\]
is an isomorphism.
\end{lemma}

\begin{proof}
Since \(R(p)=1\), the polynomials \(L\) and \(Q\) are coprime. Let
\((\dot L,\dot Q)\in T_pY\) lie in the kernel of \(dM_p\).
Differentiating \(LQ\) gives
\[
\dot L\,Q+L\,\dot Q=0.
\]
Thus \(L\mid \dot LQ\). Because \(\gcd(L,Q)=1\), unique factorization in
\(\C[U,V]\) gives \(L\mid\dot L\). Both are linear, so
\(\dot L=tL\) for some \(t\in\C\). The differentiated product identity then
gives \(\dot Q=-tQ\).

The only possible infinitesimal kernel direction is therefore the gauge
direction. But differentiating the resultant scaling law at \(\lambda=1\)
gives
\[
dR_p(tL,-tQ)=tR(p)=t.
\]
Tangency to \(Y=\{R=1\}\) requires \(dR_p(\dot L,\dot Q)=0\), hence \(t=0\).
Thus the differential is injective.

The hypersurface \(Y\) is smooth: \(R\) is homogeneous of degree three, so
Euler's identity gives
\[
aR_a+bR_b+cR_c+dR_d+eR_e=3R=3
\]
on \(Y\); consequently \(\nabla R\neq0\). Hence
\(\dim T_pY=4\), equal to the target dimension. The injective differential is
therefore an isomorphism.
\end{proof}

The pipeline control has now passed: after the trivial gauge is removed,
multiplication loses no infinitesimal information. Nevertheless, the three
distinct marked states constructed above remain. The resultant normalization
removes a continuous redundancy but cannot restore the erased discrete label.

Indeed, each marked factorization of a cubic with three distinct roots has
nonzero resultant. Choosing \(\lambda=R(L,Q)^{-1}\) in the scaling law
normalizes it to \(R=1\) without changing the product, so all three marked
states survive inside \(Y\).

\section{The affine boundary and residual closure}

Let \(\mathcal V\subset\C^4\) be the affine hyperplane where the
\(U^2V\)-coefficient equals one:
\[
\mathcal V=\{(f,g,h,i)\in\C^4:g=1\}.
\]
Its normalized preimage is
\[
X
=Y\cap M^{-1}(\mathcal V)
=\left\{
(a,b,c,d,e)\in\C^5:
\begin{array}{l}
a^2e-abd+b^2c=1,\\
ad+bc=1
\end{array}
\right\}.
\]

The double-root slice enters here through the middle coefficient \(ad+bc\).
At the boundary \(a=0\), the two constraints become
\[
b^2c=1,\qquad bc=1,
\]
and therefore
\[
b=c=1.
\]
The generic cubic-quadratic intersection does not branch in the affine
boundary: it collapses to a single base point in \((b,c)\), with \(d,e\) free.
This is the closure seam.

The first two normalized residuals are
\[
b-1=ay,
\qquad
c-1+\frac32ay=a^2z.
\]
Promoting \(y\) and \(z\) to coordinates yields the following global chart.

\begin{lemma}[Polynomial closure of the slice]\label{lem:closure}
The map \(\Phi:\C^3\to X\), defined by
\[
a=a,\qquad
b=1+ay,\qquad
c=1-\frac32ay+a^2z,
\]
\[
d=\frac12y-az+\frac32ay^2-a^2yz,
\]
\[
e=-2z+4y^2-4ayz+3ay^3-2a^2y^2z,
\]
is a polynomial isomorphism.
\end{lemma}

\begin{proof}
Substitution in the displayed formulas gives the two polynomial identities
\[
ad+bc=1,
\qquad
2bd-ae=y.
\]
Because \(b=1+ay\), these identities imply
\begin{align*}
a^2e-abd+b^2c
  &=a^2e-abd+b(1-ad)\\
  &=a^2e-2abd+b\\
  &=b-ay\\
  &=1.
\end{align*}
Thus \(\Phi(\C^3)\subseteq X\).

Conversely, take \((a,b,c,d,e)\in X\) and define polynomially
\[
y=2bd-ae,
\qquad
z=2d^2+3dy+e\left(c-\frac32\right).
\]
Using the two equations of \(X\),
\begin{align*}
ay
  &=2abd-a^2e\\
  &=abd+b^2c-1\\
  &=b(ad+bc)-1\\
  &=b-1.
\end{align*}
Hence \(b=1+ay\). Next,
\begin{align*}
\frac32y-d-yc
  &=d(3b-1-2bc)+ae\left(c-\frac32\right)\\
  &=a\left(3dy+2d^2+e\left(c-\frac32\right)\right)\\
  &=az,
\end{align*}
where we used \(bc=1-ad\) and \(b-1=ay\). Since \(b=1+ay\) and
\(bc=1-ad\),
\[
c+ayc=1-ad,
\qquad\text{so}\qquad
c-1=-a(d+yc).
\]
Together with the preceding identity, this gives
\[
c-1+\frac32ay=a^2z,
\]
which recovers the displayed formula for \(c\).

If \(a\neq0\), substituting the recovered \(b,c\) into \(ad+bc=1\) gives the
displayed formula for \(d\). The relation \(y=2bd-ae\) then gives the displayed
formula for \(e\).

It remains to treat the boundary rather than cancel \(a\). If \(a=0\), the
boundary equations give \(b=c=1\). The polynomial inverse coordinates then
give
\[
y=2d,
\qquad
z=8d^2-\frac12e.
\]
Therefore
\[
d=\frac12y,
\qquad
e=-2z+4y^2,
\]
which are exactly the \(a=0\) cases of the displayed formulas. Thus the two
polynomial coordinates in the inverse construction define an inverse to
\(\Phi\) on all of \(X\), including the degenerate fiber.
\end{proof}

This is the differential-closure step of the pipeline: the first residual is
divisible by \(a\), the corrected second residual is divisible by \(a^2\), and
the resulting coordinates close after two levels.

\section{The induced map and the structural Jacobian}

Identify \(\mathcal V\) with \(\C^3\) by deleting its fixed second coordinate.
The induced map
\[
P:\C^3\longrightarrow\C^3
\]
is
\[
P(a,y,z)=\bigl(G_1(a,y,z),G_2(a,y,z),G_3(a,y,z)\bigr),
\]
where
\[
G_1=a-\frac32a^2y+a^3z,
\]
\[
G_2=\frac12y-3az+6ay^2-6a^2yz
     +\frac92a^2y^3-3a^3y^2z,
\]
\[
G_3=-2z+4y^2-6ayz+7ay^3-6a^2y^2z
     +3a^2y^4-2a^3y^3z.
\]
These are simply the remaining coefficients
\[
(G_1,G_2,G_3)=(ac,\ ae+bd,\ be)
\]
after substitution of the polynomial chart.

By Lemma~\ref{lem:local-fidelity}, \(M|_Y\) is locally invertible. Restricting
it to the inverse image of the smooth hyperplane \(\mathcal V\) preserves local
invertibility. Lemma~\ref{lem:closure} supplies polynomial coordinates on that
inverse image. Therefore \(DP\) is invertible at every point of \(\C^3\), so
\(\det DP\) is a polynomial with no complex zero.

A nonconstant polynomial over \(\C\) has a zero after restriction to a suitable
complex line. Hence every nowhere-zero polynomial on \(\C^3\) is constant. At
the origin, the displayed linear terms give
\[
DP(0,0,0)
=
\begin{pmatrix}
1&0&0\\
0&\frac12&0\\
0&0&-2
\end{pmatrix}.
\]
Consequently,
\[
\det DP\equiv -1.
\]

The large cancellation in the expanded Jacobian is therefore not accidental:
it is forced by local fidelity, polynomial closure, and the fact that a
nowhere-zero complex polynomial is a unit.

\section{Compact certificate}

Introduce source coordinates \((x,y,t)\) by
\[
(a,y,z)=\left(x,y,-\frac12t\right)
\]
and apply the linear target change
\[
S(u,v,w)=(w,2v,2u).
\]
Define
\[
F=S\circ P\circ T,
\qquad
T(x,y,t)=\left(x,y,-\frac12t\right).
\]
Expanding gives the compact certificate
\[
F_1(x,y,t)
=(1+xy)^3t+y^2(1+xy)(4+3xy),
\]
\[
F_2(x,y,t)
=y+3x(1+xy)^2t+3xy^2(4+3xy),
\]
\[
F_3(x,y,t)
=2x-3x^2y-x^3t.
\]

Since
\[
\det DS=-4,
\qquad
\det DP=-1,
\qquad
\det DT=-\frac12,
\]
the chain rule gives
\[
\det DF=(-4)(-1)\left(-\frac12\right)=-2.
\]

It remains to expose the global alias. Let
\[
p_0=\left(0,0,-\frac14\right),
\qquad
p_\varepsilon=
\left(\varepsilon,-\frac32\varepsilon,\frac{13}{2}\right),
\quad \varepsilon\in\{1,-1\}.
\]
At \(p_0\), direct substitution gives
\(F(p_0)=(-\tfrac14,0,0)\). At \(p_\varepsilon\),
\[
xy=-\frac32,\qquad 1+xy=-\frac12,\qquad y^2=\frac94.
\]
Therefore
\[
F_1(p_\varepsilon)
=-\frac{13}{16}+\frac{9}{16}
=-\frac14,
\]
\[
F_2(p_\varepsilon)
=-\frac{3\varepsilon}{2}
 +\frac{39\varepsilon}{8}
-\frac{27\varepsilon}{8}
=0,
\]
\[
F_3(p_\varepsilon)
=2\varepsilon+\frac{9\varepsilon}{2}
-\frac{13\varepsilon}{2}
=0.
\]
Thus the three distinct points satisfy
\[
F(p_0)=F(p_1)=F(p_{-1})
=\left(-\frac14,0,0\right).
\]

The map \(F\) is polynomial and has constant nonzero Jacobian determinant, yet
it is not injective. This proves Theorem~\ref{thm:counterexample}. For \(n>3\),
the product map
\[
(F,\id_{\C^{\,n-3}}):\C^n\longrightarrow\C^n
\]
has the same collision and the same nonzero constant Jacobian up to the
identity factor.

\section{Closure audit}

The \UPSN{} requires closure before presentation. The proof satisfies the
audit.

\begin{auditbox}
\begin{itemize}
  \item \textbf{Exact claim:} an explicit counterexample over \(\C^3\) was
        constructed.
  \item \textbf{Domain:} every displayed map is polynomial on the stated
        affine space.
  \item \textbf{Local condition:} \(\det DF\equiv-2\), established
        structurally by the chain rule.
  \item \textbf{Global failure:} three distinct rational points have one exact
        image.
  \item \textbf{Gauge control:} the continuous scaling ambiguity was removed
        before the discrete collision was used.
  \item \textbf{Boundary case:} \(a=0\) was handled separately; no illegal
        division by \(a\) occurs.
  \item \textbf{Circularity:} affineness, local fidelity, and noninjectivity
        were proved by separate obligations.
  \item \textbf{Theorem/experiment separation:} symbolic computation was used
        only as an independent exact audit; the written proof does not depend
        on numerical evidence.
\end{itemize}
\end{auditbox}

\section{Interpretation}

The counterexample has a compact conceptual form.

\begin{mechanismbox}
\centering\bfseries\color{JGPInk}
Local differential fidelity does not imply global state fidelity.
\end{mechanismbox}

The Jacobian records whether the visible representation collapses
infinitesimal directions. It cannot record a discrete state label that the
representation has already erased. Here that label is the choice of which root
belongs to the marked linear factor.

The resultant normalization proves that no continuous ambiguity remains. The
affine slice proves that the constrained state space is still ordinary
three-dimensional affine space. The collision proves that the erased discrete
label remains globally consequential. The final polynomial is the compressed
certificate of that mechanism.

\section*{Provenance and disclosure}
\addcontentsline{toc}{section}{Provenance and disclosure}

The explicit degree-seven counterexample displayed above was publicly
attributed to the Fable AI system in July 2026 and was subsequently discussed
in geometric form by Terence Tao and David Speyer
\cite{Tao2026,Speyer2026}. This note does
\textbf{not} claim discovery of the displayed map. It gives an independently
organized proof through Jon Poplett's established pipeline.

\begin{routebox}
\centering\sffamily\bfseries\color{JGPDarkBlue}
representation \(\longrightarrow\) constraint
\(\longrightarrow\) local control
\(\longrightarrow\) collision\\[2pt]
\(\longrightarrow\) independent witness
\(\longrightarrow\) boundary closure
\(\longrightarrow\) exact certificate
\end{routebox}

\begin{itemize}
  \item \textbf{Primary author and investigator:} Jon Poplett.
  \item \textbf{Assistance:} AI-assisted drafting, algebra checking, and
        document preparation.
  \item \textbf{Verification:} all decisive polynomial identities were checked
        exactly over the rational numbers; no floating-point inference enters
        the theorem.
\end{itemize}

\addcontentsline{toc}{section}{References}

\begin{thebibliography}{9}

\bibitem{Tao2026}
Terence Tao,
\emph{A digestion of the Jacobian conjecture counterexample},
What's New, 21 July 2026.
\href{https://terrytao.wordpress.com/2026/07/21/a-digestion-of-the-jacobian-conjecture-counterexample/}{Online article}.

\bibitem{Speyer2026}
David E. Speyer,
\emph{The new counterexample to the Jacobian conjecture},
Secret Blogging Seminar, 20 July 2026.
\href{https://sbseminar.wordpress.com/2026/07/20/the-new-counterexample-to-the-jacobian-conjecture/}{Online article}.

\bibitem{PoplettProtocol2026}
Jon Poplett,
\emph{Universal Proof Selection Network---Proof Construction Protocol},
unpublished working protocol, July 2026.

\bibitem{PoplettCheckpoint2026}
Jon Poplett,
\emph{Checkpoint Process Reports 01--11: Finite Local Neighborhood Aliasing},
unpublished research notes, 13 June 2026.

\end{thebibliography}

\end{document}
